%% ============================================================
%%  8.  RELATED WORK
%% ============================================================
\section{Related Work}
\label{sec:related}

\paragraph{Multi-agent LLM debate.}
\citet{du2023improving} demonstrated that iterative debate among multiple
LLM instances improves factual accuracy and reasoning quality over
single-agent chains.
\citet{liang2023encouraging} showed that divergent thinking emerges when
agents are explicitly prompted to disagree.
\citet{chan2023chateval} applied multi-agent debate to evaluation tasks.
\citet{xiong2023examining} studied inter-agent consistency and convergence
dynamics.
\citet{chen2024reconcile} introduced round-table conferencing across diverse
LLMs.
A key limitation across this prior work is its focus on \emph{outcome}
optimisation: whether the final answer is correct or preferred.
Our closest predecessor, \citeauthor{diagnostic2026}~(\citeyear{diagnostic2026}),
shifts to \emph{process} evaluation grounded in deliberative axioms, but
stops at measurement.
We extend this to adaptive control.

\paragraph{Deliberative democracy and social choice.}
Deliberative polling~\cite{fishkin2009deliberation} studies how structured
small-group discussion changes individual opinions toward more informed
positions—an analogue to our debate protocol.
\citet{arrow1950difficulty} and \citet{sen1970collective} establish axiomatic
foundations for collective decision-making that motivate our eight deliberative
axioms (Table~\ref{tab:axioms}).
\citet{irving2018ai} propose adversarial debate as a scalable AI safety
mechanism, arguing that a human judge can identify the correct answer when
two agents argue opposing sides.
Our loser-resistance mechanism (Exp2a/2b) is inspired by this adversarial
framing but conditions activation on process diagnostics to avoid the
indiscriminate entrenchment that hurts PER in Exp2a.

\paragraph{Adaptive intervention and selective reasoning.}
\citet{lightman2023verify} demonstrate that process-level verification
(step-by-step checking) outperforms outcome-level verification in
mathematical reasoning.
\citet{wang2023selfconsistency} show that sampling multiple reasoning paths
and selecting by majority vote improves reliability—our baseline debate
protocol is related but adds iterative interaction.
\citet{klein2007performing} introduces the pre-mortem technique—imagining
that a decision has failed and reasoning backward to identify causes—which
directly motivates our devil's advocate design.
To our knowledge, no prior work uses live process diagnostics as adaptive
control signals within a running multi-agent debate.
